\documentclass[french]{article}
\usepackage[utf8]{inputenc}
\usepackage[T1]{fontenc}
\usepackage{lmodern}
\usepackage{geometry}
\usepackage{graphicx}
\usepackage{babel}
\usepackage{amsmath}
\usepackage{amsthm}
\usepackage{amssymb}
\usepackage{hyperref}
\usepackage{stmaryrd}
\usepackage{enumitem}
\usepackage{tcolorbox}
\usepackage{dsfont}
\usepackage{multirow}
\usepackage{etoc}
\usepackage{titlesec}
\usepackage{esint}
\usepackage{cancel}
\usepackage{mathdots}

\setlist[itemize]{label=$\bullet$}


\renewcommand{\P}{\mathbb{P}}
\newcommand{\N}{\mathbb{N}}
\newcommand{\Z}{\mathbb{Z}}
\newcommand{\Q}{\mathbb{Q}}
\newcommand{\R}{\mathbb{R}}
\newcommand{\C}{\mathbb{C}}
\newcommand{\K}{\mathbb{K}}
\renewcommand{\L}{\mathbb{L}}
\newcommand{\U}{\mathbb{U}}
\newcommand{\st}{^*}
\newcommand{\tq}{\middle|}
\newcommand{\inv}{^{-1}}
\newcommand{\E}{\mathbb{E}}
\newcommand{\F}{\mathbb{F}}
\newcommand{\V}{\mathbb{V}}
\renewcommand{\ker}{\text{Ker}}
\newcommand{\im}{\text{Im}}
\newcommand{\card}{\text{Card}}
\newcommand{\Tr}{\text{Tr}}

\title{TD Euclidiens\\ Exercice 47}
\date{}
\begin{document}
\maketitle

\section{Énoncé $\Delta\star$: Égalités de Courant-Fischer}
On considère $u\in S(E)$ avec $E$ euclidien de dimension $n$. On note $\lambda_1\leqslant\ldots\leqslant\lambda_n$ ses valeurs propres. On note $\mathcal{F}_k$ l'ensemble des sous-espaces vectoriels de $E$ de dimension $k$ et $S$ la sphère unité de $E$. Montrer que :  $$\lambda_k=\inf_{F\in\mathcal{F}_k}\sup_{x\in S\cap F}(x|u(x))$$ puis montrer que $$\lambda_{n+1-k}=\sup_{F\in\mathcal{F}_k}\inf_{x\in S\cap F}(x|u(x))$$
\section{Solution}
Pour cela, prendre une BON de diagonalisation de $u$ et montrer que $S\cap F$ a une intersection avec $\text{Vect}(e_k,\ldots,e_n)$ non nulle (Grassmann) puis prendre un vecteur de norme $1$ dedans, ce qui montre que tous les sup sont atteints. Pour le sens réciproque, s'intéresser à $\text{Vect}(e_1,\ldots,e_k)$ et faire le produit scalaire (comme dans le cours). Pour montrer la deuxième égalité : $$\lambda_{n+1-k}=\sup_{F\in \mathcal{F}_k}\inf_{x\in S\cap F}(x|u(x))$$ on peut appliquer la première égalité à $-u$, ce qui échange essentiellement les $\sup$ et les $\inf$.
\end{document}
